% GENERATED FILE. Edit canonical lessons, then run npm run generate:books.
% canonical-source-hash: fnv1a64:70c63052f2028df1
% canonical-lessons: TA-C467-kintal, TA-C467-cettai, TA-C467-naccari, TA-C467-ularu, TA-C467-kurukuru

\chapter{Teasing and Pestering}
\label{ch:ta-teasing-and-pestering}
\hlchaptermodality{467}

\begin{chapteropening}
\textbf{By the end of this chapter:} \emph{I can talk about teasing, pranks, nagging and talking nonsense.}

Complete the last lesson of chapter 467: I can talk about teasing, pranks, nagging and talking nonsense.
\end{chapteropening}

\section[\texorpdfstring{\ta{கிண்டல்} (kiṇṭal)}{கிண்டல் (kiṇṭal)}]{\ta{கிண்டல்} (kiṇṭal) — teasing; mockery}

\label{lesson:TA-C467-kintal}



\begin{quote}
Before the new one: say the Tamil for to cower, then the Tamil for to threaten.
\end{quote}

\subsection*{You'll want to know: \ta{கிண்டல்}}
\textbf{\ta{கிண்டல்}} — \emph{kiṇṭal} — \textquotedblleft{}teasing; mockery\textquotedblright{}.

\begin{quote}
\ta{என்} \ta{தம்பி} \ta{என்} \ta{சமையலைக்} \ta{கிண்டல்} \ta{செய்தான்}.
\end{quote}

\begin{quote}
\emph{eṉ tampi eṉ camaiyalaik kiṇṭal ceytāṉ.}
\end{quote}

\begin{quote}
My younger brother made fun of my cooking.
\end{quote}

\ta{கிண்டல்} \ta{செய்}, to tease.

\begin{cousinweb}
\textbf{\ta{கிண்டல்}} (\emph{kiṇṭal}) is native Dravidian, from \ta{கிண்டு} \emph{kiṇṭu}, \textquotedblleft{}to poke, to stir\textquotedblright{} (DEDR 1542); the entry gives \ta{கிண்டல்} as stirring up mischief.
\end{cousinweb}

\begin{grammarlens}[title={What you've built}]
One more word for describing what you lived and felt.
\end{grammarlens}

\subsection*{Your turn}
\begin{itemize}
\raggedright
  \item \emph{Say it:} \emph{kiṇṭal}
  \item \emph{Say it:} \emph{kiṇṭal}, once more
  \item \emph{Say it:} \emph{eṉ tampi eṉ camaiyalaik kiṇṭal ceytāṉ.} — My younger brother made fun of my cooking.
  \item \emph{Recall:} say the Tamil for to cower, then the Tamil for to threaten, then say \emph{kiṇṭal} again
\end{itemize}

\begin{tcolorbox}[breakable,colback=teal!4,colframe=teal!35!black,title={Before you move on}]
What does \ta{கிண்டல்} mean? (\textquotedblleft{}Teasing; mockery\textquotedblright{}.) Say it once more.
\end{tcolorbox}
\section[\texorpdfstring{\ta{சேட்டை} (cēṭṭai)}{சேட்டை (cēṭṭai)}]{\ta{சேட்டை} (cēṭṭai) — naughtiness; pranks}

\label{lesson:TA-C467-cettai}



\begin{quote}
Before the new one: say the Tamil for to threaten, then the Tamil for teasing.
\end{quote}

\subsection*{You'll want to know: \ta{சேட்டை}}
\textbf{\ta{சேட்டை}} — \emph{cēṭṭai} — \textquotedblleft{}naughtiness; pranks\textquotedblright{}.

\begin{quote}
\ta{வகுப்பில்} \ta{அந்தப்} \ta{பையன்} \ta{சேட்டை} \ta{செய்தான்}.
\end{quote}

\begin{quote}
\emph{vakuppil antap paiyaṉ cēṭṭai ceytāṉ.}
\end{quote}

\begin{quote}
That boy played pranks in class.
\end{quote}

\ta{சேட்டை} \ta{செய்}, to be naughty.

\begin{cousinweb}
\textbf{\ta{சேட்டை}} (\emph{cēṭṭai}) is borrowed from Sanskrit \emph{cēṣṭā}, \textquotedblleft{}action, behaviour\textquotedblright{} (CDIAL 4913); Tamil narrows it to mischief.
\end{cousinweb}

\begin{grammarlens}[title={What you've built}]
One more word for describing what you lived and felt.
\end{grammarlens}

\subsection*{Your turn}
\begin{itemize}
\raggedright
  \item \emph{Say it:} \emph{cēṭṭai}
  \item \emph{Say it:} \emph{cēṭṭai}, once more
  \item \emph{Say it:} \emph{vakuppil antap paiyaṉ cēṭṭai ceytāṉ.} — That boy played pranks in class.
  \item \emph{Recall:} say the Tamil for to threaten, then the Tamil for teasing, then say \emph{cēṭṭai} again
\end{itemize}

\begin{tcolorbox}[breakable,colback=teal!4,colframe=teal!35!black,title={Before you move on}]
What does \ta{சேட்டை} mean? (\textquotedblleft{}Naughtiness; pranks\textquotedblright{}.) Say it once more.
\end{tcolorbox}
\section[\texorpdfstring{\ta{நச்சரி} (naccari)}{நச்சரி (naccari)}]{\ta{நச்சரி} (naccari) — to nag; to pester}

\label{lesson:TA-C467-naccari}



\begin{quote}
Before the new one: say the Tamil for teasing, then the Tamil for naughtiness.
\end{quote}

\subsection*{You'll want to know: \ta{நச்சரி}}
\textbf{\ta{நச்சரி}} — \emph{naccari} — \textquotedblleft{}to nag; to pester\textquotedblright{}.

\begin{quote}
\ta{குழந்தை} \ta{பொம்மை} \ta{வாங்கச்} \ta{சொல்லி} \ta{அம்மாவை} \ta{நச்சரித்தது}.
\end{quote}

\begin{quote}
\emph{kuḻantai pommai vāṅkac colli ammāvai naccarittatu.}
\end{quote}

\begin{quote}
The child pestered its mother to buy a toy.
\end{quote}

Past (I …): \ta{நச்சரித்தேன்}, \emph{naccarittēṉ}.

The one pestered takes -\ta{ஐ}.

\begin{cousinweb}
\textbf{\ta{நச்சரி}} (\emph{naccari}) is native Tamil, usually traced to \ta{நச்சு} \emph{naccu}, \textquotedblleft{}to tease, to vex\textquotedblright{} (DEDR 3577), which does not list this form itself.
\end{cousinweb}

\begin{grammarlens}[title={What you've built}]
One more word for describing what you lived and felt.
\end{grammarlens}

\subsection*{Your turn}
\begin{itemize}
\raggedright
  \item \emph{Say it:} \emph{naccari}
  \item \emph{Say it:} \emph{naccari}, once more
  \item \emph{Say it:} \emph{kuḻantai pommai vāṅkac colli ammāvai naccarittatu.} — The child pestered its mother to buy a toy.
  \item \emph{Recall:} say the Tamil for teasing, then the Tamil for naughtiness, then say \emph{naccari} again
\end{itemize}

\begin{tcolorbox}[breakable,colback=teal!4,colframe=teal!35!black,title={Before you move on}]
What does \ta{நச்சரி} mean? (\textquotedblleft{}To nag; to pester\textquotedblright{}.) Say it once more.
\end{tcolorbox}
\section[\texorpdfstring{\ta{உளறு} (uḷaṟu)}{உளறு (uḷaṟu)}]{\ta{உளறு} (uḷaṟu) — to babble; to talk nonsense}

\label{lesson:TA-C467-ularu}



\begin{quote}
Before the new one: say the Tamil for naughtiness, then the Tamil for to nag.
\end{quote}

\subsection*{You'll want to know: \ta{உளறு}}
\textbf{\ta{உளறு}} — \emph{uḷaṟu} — \textquotedblleft{}to babble; to talk nonsense\textquotedblright{}.

\begin{quote}
\ta{காய்ச்சலில்} \ta{அவன்} \ta{ஏதோ} \ta{உளறினான்}.
\end{quote}

\begin{quote}
\emph{kāyccalil avaṉ ētō uḷaṟiṉāṉ.}
\end{quote}

\begin{quote}
In his fever he babbled something.
\end{quote}

Past (I …): \ta{உளறினேன்}, \emph{uḷaṟiṉēṉ}.

Also to blurt out: \ta{ரகசியத்தை} \ta{உளறினான்}, he let the secret slip.

\begin{cousinweb}
\textbf{\ta{உளறு}} (\emph{uḷaṟu}) is native Dravidian (DEDR 760), in an entry of words for howling and senseless noise.
\end{cousinweb}

\begin{grammarlens}[title={What you've built}]
One more word for describing what you lived and felt.
\end{grammarlens}

\subsection*{Your turn}
\begin{itemize}
\raggedright
  \item \emph{Say it:} \emph{uḷaṟu}
  \item \emph{Say it:} \emph{uḷaṟu}, once more
  \item \emph{Say it:} \emph{kāyccalil avaṉ ētō uḷaṟiṉāṉ.} — In his fever he babbled something.
  \item \emph{Recall:} say the Tamil for naughtiness, then the Tamil for to nag, then say \emph{uḷaṟu} again
\end{itemize}

\begin{tcolorbox}[breakable,colback=teal!4,colframe=teal!35!black,title={Before you move on}]
What does \ta{உளறு} mean? (\textquotedblleft{}To babble; to talk nonsense\textquotedblright{}.) Say it once more.
\end{tcolorbox}
\section[\texorpdfstring{\ta{குறுகுறு} (kuṟukuṟu)}{குறுகுறு (kuṟukuṟu)}]{\ta{குறுகுறு} (kuṟukuṟu) — to tickle; to prick}

\label{lesson:TA-C467-kurukuru}



\begin{quote}
Before the new one: say the Tamil for to nag, then the Tamil for to babble.
\end{quote}

\subsection*{You'll want to know: \ta{குறுகுறு}}
\textbf{\ta{குறுகுறு}} — \emph{kuṟukuṟu} — \textquotedblleft{}to tickle; to prick\textquotedblright{}.

\begin{quote}
\ta{பொய்} \ta{சொன்னதால்} \ta{என்} \ta{மனம்} \ta{குறுகுறுத்தது}.
\end{quote}

\begin{quote}
\emph{poy coṉṉatāl eṉ maṉam kuṟukuṟuttatu.}
\end{quote}

\begin{quote}
Because I had told a lie, my conscience pricked me.
\end{quote}

Past (it …): \ta{குறுகுறுத்தது}, \emph{kuṟukuṟuttatu}.

A doubled sound word; the subject is the mind or a part of the body, so the past is \ta{குறுகுறுத்தது}.

\begin{cousinweb}
\textbf{\ta{குறுகுறு}} (\emph{kuṟukuṟu}) is a doubled native word whose origin is not settled; DEDR has no entry for this sense.
\end{cousinweb}

\begin{grammarlens}[title={What you've built}]
One more word for describing what you lived and felt.
\end{grammarlens}

\subsection*{Your turn}
\begin{itemize}
\raggedright
  \item \emph{Say it:} \emph{kuṟukuṟu}
  \item \emph{Say it:} \emph{kuṟukuṟu}, once more
  \item \emph{Say it:} \emph{poy coṉṉatāl eṉ maṉam kuṟukuṟuttatu.} — Because I had told a lie, my conscience pricked me.
  \item \emph{Recall:} say the Tamil for to nag, then the Tamil for to babble, then say \emph{kuṟukuṟu} again
\end{itemize}

\begin{tcolorbox}[breakable,colback=teal!4,colframe=teal!35!black,title={Before you move on}]
What does \ta{குறுகுறு} mean? (\textquotedblleft{}To tickle; to prick\textquotedblright{}.) Say it once more.
\end{tcolorbox}
